\chapter{The filter}\label{ch:filter}
\begin{goals}
\textbf{In this chapter} you narrow the table to the rows you care about: typing a filter, building one
from the grid, editing it as controls, reading its errors, and knowing what it costs on a large file.
\end{goals}

The filter is a condition on a row. lamina keeps the rows for which it is true, in their order in the
file, and everything else (the grid, the sort, the statistics, the instruments, the export) works on
those rows. The condition is written in lamina's expression language, which \cref{ch:expr} describes in
full; this chapter is about using it.

\section{The filter box}\label{sec:filter-box}

Type a condition in the box on the toolbar and press \key{Enter}:

\begin{codeblock}
LITO = "SULPHIDE" and Cu_pct > 0.5
\end{codeblock}

The grid shows the matching rows, and the status line counts them:
\code{5,362 of 120,000 rows (filtered)}. \key{Esc} in the box clears the filter, and so does the \ui{×}
at its end; \ui{Filter} applies it, as \key{Enter} does. The box keeps its own undo: \key{Ctrl}+\key{Z}
brings back what it held before.

As you type, the box colours the expression: column names, text, numbers, keywords and functions each
in their own colour, so a column name typed where text was meant shows at once. A list of suggestions
follows the cursor: column names, functions, keywords, and after \code{LITO =} or \code{LITO in (}, the
column's own values, already quoted. \key{↑} and \key{↓} move through it, \key{Tab} or \key{Enter}
takes one, and \key{Esc} closes the list (a second \key{Esc} then clears the filter). While the list is
open, \key{Enter} takes a suggestion rather than applying the filter.

\begin{note}
A word on its own is a column; a word in quotes is text. \code{LITO = "OXIDE"} compares the column with
the text; \code{LITO = OXIDE} looks for a column called OXIDE. Picking a value from the suggestions puts
the quotes in for you. \Cref{sec:names} has the details.
\end{note}

\section{When the filter is wrong}\label{sec:filter-errors}

A filter that can't be read turns the box red and the status line says why, prefixed \code{filter:}. The
grid keeps showing what it showed before.

\begin{center}\small
\begin{tabularx}{\linewidth}{@{}>{\ttfamily}p{0.38\linewidth}L@{}}
\toprule
\normalfont\sffamily\bfseries you typed & \sffamily\bfseries lamina says \\
\midrule
cu\_pc > 1 & \code{unknown column: cu\_pc — did you mean Cu\_pct?} \\
Cu\_pct > & \code{expr parse error: unexpected end of expression} \\
select * from t where Cu\_pct > 1 & \code{the filter IS the WHERE clause — drop the rest, just: Cu\_pct > 1} \\
ZONE matches "(a+)+" & a refusal of the pattern, which could take forever on some values
(\cref{sec:text-match}) \\
\bottomrule
\end{tabularx}
\end{center}

\begin{caution}
For a value typed without quotes, the suggestion can mislead: \code{LITO = OXIDE} gives \code{unknown
column: OXIDE — did you mean ID?}, because OXIDE is two letters from no column and three from \col{ID}.
What it means is that \code{OXIDE} needs quotes.
\end{caution}

\section{Filters from the grid}\label{sec:filter-from-grid}

You don't have to type a filter. Several things in the grid write one into the box and apply it:

\begin{description}
\item[A cell's menu] \ui{Filter LITO = OXIDE}: the rows with that cell's value. lamina writes it in
the C-style spelling, \code{LITO == "OXIDE"}, which means the same as \code{=}.
\item[A header's menu] \ui{Filter by Cu\_pct…} puts the column's name in the box, ready for you to type
the rest.
\item[A distribution] a click or a drag over the glyph above a column (\cref{sec:gutter-use}).
\item[The statistics window] of a text column: click values in its list to build \code{LITO in "OXIDE",
"TRANSITION"}, one click per value (\cref{sec:one-column}).
\item[Group by] a click on a group filters to it (\cref{ch:groupby}).
\end{description}

Each of these replaces the filter in the box; none adds to it. To combine, edit the box.

\section{The filter as controls}\label{sec:drawer}

The \ui{▾} at the end of the filter box opens the filter's \term{drawer}: the current expression drawn as
controls, one for each condition. A range of numbers (\code{Cu\_pct between 0.5 and 2},
\code{Z >= 450}) becomes a slider over the column's range; a set of values (\code{LITO in (…)}) a
row of value chips over the column's values, selected or not. Each condition has its operator in a list
beside it (change \ui{between} to \ui{>=}, \ui{in} to \ui{not in}) and \ui{×} to remove it; \ui{add
condition…} adds one for another column, joined by the \ui{and} or \ui{or} beside it. The expression itself
is shown above the controls and can be edited there. Moving a control rewrites that part of the expression,
and the rest is left as you wrote it, comments and all. \ui{reset} goes back to the filter the drawer
opened with.

\shot[0.75\linewidth]{drawer}{The filter drawer for \code{LITO in ("SULPHIDE", "TRANSITION") and
Cu\_pct between 0.5 and 2}: checkboxes for the lithologies, a slider for copper.}

The sliders' ranges come from the sample behind the distributions (\cref{sec:gutter-sample}), so they
are approximate; the numbers they write into the expression are exact. The filter is applied a moment
after a control stops moving, not at every step of a drag, since each application is a scan of the file.

\section{What a filter scans}\label{sec:filter-scan}

A filter is a scan: it reads the file once and keeps the location of each matching row
(\cref{sec:scans}). On a large file the status line counts it, \code{filtering… 42\% · Esc to cancel}, and
\key{Esc} stops it, leaving the previous view.

How much of the file a filter reads depends on the format:

\begin{itemize}
\item On delimited text, every row is read and split, whatever the filter names.
\item On a \file{.dm}, only the fields the filter names are read (\cref{sec:strided}).
\item On Parquet, only those columns are read, and row groups whose statistics rule a match out are
skipped (\cref{sec:pruning}).
\end{itemize}

With a calculated column defined, a \file{.dm} or Parquet filter reads every column
(\cref{ch:calc}).

\section{Filter and sort together}

A filter and a sort compose. The sort orders the matches, and changing the filter sorts the new matches
again. On a file too large to sort whole (\cref{sec:huge-limits}), filter first and sort the result:
filtering keeps 24 bytes per match, and the sort then has only the matches to order.

\begin{tryit}
Three ways to the same rows.
\begin{enumerate}
\item In the sample, type \code{LITO = "TRANSITION" and Au\_gpt > 1} and press \key{Enter}. Note the
count in the status line.
\item Press \key{Esc}. Right-click a TRANSITION cell in \col{LITO}, choose \ui{Filter LITO =
TRANSITION}, then add \code{and Au\_gpt > 1} to the end of what the box holds and press \key{Enter}.
The same count.
\item Open the drawer with \ui{▾}. Drag the \col{Au\_gpt} slider's lower end, and watch the count follow.
\end{enumerate}
\end{tryit}
